% TA-IND-02 -- Informe Tecnico Individual: Analisis de Consistencia y Replicacion en BDD Distribuidas
% Estudiante: Jeremy Alexis Alvarez Parraga -- Equipo PFC: ACC (Soporte Tecnico ISP)
% Repositorio: https://github.com/DeJere/TA-IND-02-Informe-T-cnico-Individual-.git
\documentclass[11pt,a4paper]{article}

\usepackage[utf8]{inputenc}
\usepackage[T1]{fontenc}
\usepackage[spanish,es-noquoting]{babel}
\usepackage[margin=2.5cm]{geometry}
\usepackage{lmodern}
\usepackage{microtype}
\usepackage{booktabs}
\usepackage{array}
\usepackage{caption}
\usepackage{tikz}
\usetikzlibrary{positioning,arrows.meta,fit,backgrounds,shapes.geometric}
\usepackage{enumitem}
\usepackage{setspace}

\addto\captionsspanish{\renewcommand{\tablename}{Tabla}}
\addto\captionsspanish{\renewcommand{\figurename}{Figura}}

\setstretch{1.12}
\setlength{\parskip}{0.35em}

\title{}
\date{}

\begin{document}
\shorthandoff{"}
\sloppy

% ============================================================
% PORTADA
% ============================================================
\begin{titlepage}
\centering
\vspace*{1cm}
{\large UNIVERSIDAD T\'ECNICA ESTATAL DE QUEVEDO\par}
{\normalsize Facultad de Ciencias de la Computaci\'on\par}
{\normalsize Carrera de Ingenier\'ia de Software\par}
\vspace{0.3cm}
{\normalsize Aplicaciones Distribuidas (ISR-701) -- Unidad 3: Bases de Datos Distribuidas\par}

\vspace{2.5cm}
{\Large \bfseries Informe T\'ecnico Individual\par}
\vspace{0.4cm}
{\Large \bfseries An\'alisis de Consistencia y Replicaci\'on en Bases de Datos Distribuidas\par}
\vspace{0.6cm}
{\large ¿Es el modelo de consistencia elegido el \'optimo para el dominio del sistema?\par}

\vspace{2.2cm}
\begin{tabular}{ll}
\textbf{C\'odigo:} & TA-IND-02 \\[0.3em]
\textbf{Tipo:} & Trabajo Aut\'onomo -- Individual -- Sumativa \\[0.3em]
\textbf{Unidad:} & 3 -- Bases de Datos Distribuidas \\[0.3em]
\textbf{Equipo PFC analizado:} & ACC -- Soporte T\'ecnico ISP (gesti\'on de tickets) \\[0.3em]
\textbf{Estudiante:} & Jeremy Alexis Alvarez Parraga \\[0.3em]
\textbf{Docente:} & Prof. Guerrero Ulloa G.C. \\[0.3em]
\textbf{Fecha:} & \today \\
\end{tabular}

\vspace{1.2cm}
{\small\texttt{Repositorio: https://github.com/DeJere/TA-IND-02-Informe-T-cnico-Individual-.git}\par}

\vfill
{\normalsize Periodo 2026--2027\par}
\end{titlepage}

% ============================================================
% RESUMEN
% ============================================================
\section*{Resumen}

El equipo ACC del Proyecto Fin de Curso (PFC) implementa un sistema de gesti\'on de tickets de
soporte t\'ecnico para un proveedor de Internet (ISP) organizado en tres zonas operativas. Este
informe caracteriza el modelo de consistencia realmente implementado en dicho sistema: un
n\'ucleo transaccional sobre CockroachDB, particionado geogr\'aficamente por zona y replicado de
forma s\'incrona mediante consenso Raft bajo aislamiento serializable, combinado con una
coreograf\'ia de eventos as\'incrona sobre Apache Kafka entre \texttt{ticket-service} y
\texttt{ai-service} que introduce consistencia eventual en el enriquecimiento de clasificaci\'on
autom\'atica de tickets. Se ubica esta arquitectura h\'ibrida en el espectro de consistencia de
Viotti y Vukoli\'c, se interpreta la decisi\'on de dise\~no bajo los marcos CAP y PACELC, y se
contrasta el modelo actual con dos alternativas -- r\'eplica primario--secundario as\'incrona y
almacenamiento sin l\'ider por quorum -- mediante cinco criterios t\'ecnicos explicitados en una
tabla comparativa. La postura sostenida es que el modelo h\'ibrido resulta adecuado para el
dominio porque separa la consistencia fuerte que exigen los invariantes cr\'iticos del negocio
(asignaci\'on de t\'ecnicos, cumplimiento de SLA) de la consistencia relajada que es tolerable en
una tarea de enriquecimiento no bloqueante, si bien persisten limitaciones de latencia de
escritura entre zonas y de observabilidad del retraso de la coreograf\'ia.

% ============================================================
% INTRODUCCION
% ============================================================
\section{Introducci\'on}

El Proyecto Fin de Curso del equipo ACC aborda el dominio de \emph{Soporte T\'ecnico ISP}: la
gesti\'on del ciclo de vida de tickets de soporte t\'ecnico para un proveedor de Internet que
opera en tres zonas geogr\'aficas de la ciudad de Quevedo (centro, norte y sur). Cada ticket se
origina cuando un cliente reporta una incidencia (conectividad, DNS, hardware, configuraci\'on o
velocidad), se le asigna una prioridad y un t\'ecnico de la zona correspondiente, y su resoluci\'on
queda sujeta a un acuerdo de nivel de servicio (SLA) cuyo incumplimiento debe poder detectarse de
forma confiable. El sistema est\'a compuesto por tres microservicios (\texttt{auth-service},
\texttt{ticket-service} y \texttt{ai-service}), un cl\'uster distribuido de tres nodos de
CockroachDB y un backbone de mensajer\'ia Apache Kafka que conecta la creaci\'on de tickets con un
servicio de clasificaci\'on autom\'atica basado en IA.

Un sistema de esta naturaleza no puede tratarse como una \'unica base de datos monol\'itica: la
distribuci\'on geogr\'afica de los t\'ecnicos y clientes, la necesidad de tolerar la ca\'ida de un
nodo sin detener el servicio, y la existencia de una tarea de enriquecimiento (la clasificaci\'on
autom\'atica) que no deber\'ia bloquear el flujo principal de negocio, obligan a tomar decisiones
expl\'icitas sobre qu\'e garant\'ias de consistencia ofrecer y c\'omo replicar los datos. Estas
decisiones no son gratuitas: cada incremento en la fuerza de la consistencia impone mayor
coordinaci\'on entre r\'eplicas y, por lo tanto, mayor latencia y menor disponibilidad ante fallos
o particiones de red~\cite{viotti2016}.

El objetivo de este informe es doble. Primero, caracterizar con precisi\'on el modelo de
consistencia y la estrategia de replicaci\'on efectivamente implementados en el PFC del equipo
ACC, a partir de la configuraci\'on real del c\'luster de CockroachDB y de la arquitectura de
mensajer\'ia entre servicios. Segundo, evaluar de forma cr\'itica -- a la luz de los teoremas CAP
y PACELC y de al menos dos alternativas de dise\~no -- si dicho modelo constituye la elecci\'on
\'optima para el dominio de soporte t\'ecnico de un ISP, o si existen compromisos m\'as adecuados
que el equipo deber\'ia considerar.

% ============================================================
% MARCO TEORICO
% ============================================================
\section{Marco te\'orico}

\subsection{Modelos de consistencia}
\label{subsec:modelos}

La consistencia de r\'eplicas no es una propiedad binaria sino un espectro de garant\'ias con
sem\'antica formal precisa, ordenadas por su fuerza~\cite{viotti2016}. En el extremo m\'as fuerte,
la \textbf{linealizabilidad} (o consistencia fuerte) exige que toda operaci\'on parezca ejecutarse
de forma at\'omica en un instante \'unico dentro del tiempo real transcurrido entre su invocaci\'on
y su respuesta, de modo que cualquier lector observe siempre el efecto de la \'ultima escritura
confirmada. La \textbf{consistencia secuencial} relaja la referencia al tiempo real pero preserva
un orden total \'unico de operaciones, visible igual para todos los procesos. La
\textbf{consistencia causal} exige preservar \'unicamente el orden entre operaciones causalmente
relacionadas, permitiendo reordenar operaciones concurrentes e independientes. La
\textbf{consistencia de lectura de los propios cambios} (read-your-writes) es una garant\'ia
todav\'ia m\'as d\'ebil, centrada en la sesi\'on: un proceso siempre ve sus propias escrituras,
aunque no necesariamente las de otros procesos de forma inmediata. Finalmente, la
\textbf{consistencia eventual} solo garantiza que, en ausencia de nuevas escrituras, todas las
r\'eplicas converger\'an eventualmente al mismo estado, sin acotar cu\'anto tiempo tarda esa
convergencia.

Esta jerarqu\'ia tiene un correlato operacional directo: un modelo m\'as fuerte requiere que las
r\'eplicas coordinen cada operaci\'on (t\'ipicamente mediante consenso distribuido) antes de
confirmarla, lo cual incrementa la latencia y reduce la disponibilidad cuando alguna r\'eplica no
es alcanzable; un modelo m\'as d\'ebil permite confirmar operaciones localmente y propagar los
cambios de forma diferida, ganando disponibilidad y latencia a costa de exponer estados
transitoriamente divergentes entre r\'eplicas.

\subsection{Teoremas de compromiso: CAP y PACELC}
\label{subsec:cap-pacelc}

El teorema CAP, formalizado por Gilbert y Lynch a partir de la conjetura de Brewer, establece que
ante una partici\'on de red un sistema distribuido no puede garantizar simult\'aneamente
consistencia (C) y disponibilidad (A); debe sacrificar una de las dos mientras persista la
partici\'on, conservando siempre la tolerancia a particiones (P) como requisito irrenunciable en
cualquier sistema realmente distribuido~\cite{gilbert2002}. Sin embargo, CAP describe \'unicamente
el comportamiento del sistema durante una partici\'on, un evento relativamente infrecuente; no
dice nada sobre el compromiso que enfrenta el sistema en operaci\'on normal, que es el caso
dominante en la pr\'actica.

Ese vac\'io lo cubre la formulaci\'on PACELC de Abadi: si hay Partici\'on (P), el sistema elige
entre Disponibilidad (A) y Consistencia (C), igual que en CAP; pero en su defecto (E, \emph{else}),
en operaci\'on normal sin particiones, el sistema elige entre Latencia (L) y Consistencia
(C)~\cite{abadi2012}. PACELC captura as\'i que incluso sistemas que jam\'as sufren una partici\'on
real deben decidir permanentemente cu\'anta latencia est\'an dispuestos a pagar por mantener las
r\'eplicas sincronizadas. Un sistema PC/EC (como Spanner) prioriza consistencia tanto ante
particiones como en operaci\'on normal, aceptando mayor latencia; un sistema PA/EL (como Dynamo)
prioriza disponibilidad y baja latencia en ambos casos. La decisi\'on de dise\~no de cualquier PFC
debe leerse bajo ambos ejes, no solo bajo el eje de partici\'on de CAP.

\subsection{Estrategias de replicaci\'on}
\label{subsec:replicacion}

La estrategia de replicaci\'on determina c\'omo se propagan las escrituras entre r\'eplicas y, con
ello, condiciona directamente el modelo de consistencia alcanzable~\cite{vansteen2017}. En la
\textbf{replicaci\'on primario--respaldo} (\emph{leader--follower}) todas las escrituras pasan por
un nodo l\'ider que las propaga a los seguidores; si la propagaci\'on es s\'incrona (el l\'ider
espera confirmaci\'on de un qu\'orum de seguidores antes de responder al cliente) se obtienen
garant\'ias fuertes a costa de latencia, mientras que si es as\'incrona se gana disponibilidad y
latencia a costa de una ventana de inconsistencia y de posible p\'erdida de datos ante la ca\'ida
del l\'ider. La \textbf{replicaci\'on multi-l\'ider} admite varios nodos que aceptan escrituras de
forma concurrente, lo que exige resoluci\'on de conflictos y habitualmente relaja la consistencia.
La \textbf{replicaci\'on sin l\'ider por quorum} elimina la figura de l\'ider fijo: cualquier
r\'eplica puede aceptar escrituras, y la consistencia percibida depende de los par\'ametros de
qu\'orum de lectura y escritura ($R$ y $W$) frente al total de r\'eplicas $N$.

Los sistemas de referencia industrial ilustran los extremos de este espectro. Dynamo, de Amazon,
prioriza la disponibilidad mediante replicaci\'on sin l\'ider con qu\'orum ajustable y resoluci\'on
de conflictos por vectores de versi\'on, aceptando consistencia eventual como consecuencia directa
de esa elecci\'on~\cite{decandia2007}. Spanner, de Google, se ubica en el extremo opuesto: ofrece
consistencia externa (transacciones globalmente consistentes y linealizables) apoy\'andose en
consenso Paxos por fragmento de datos y en sincronizaci\'on temporal de alta precisi\'on
(TrueTime) para ordenar transacciones distribuidas sin coordinaci\'on adicional~\cite{corbett2013}.
CockroachDB, la base de datos empleada en el PFC del equipo ACC, se inspira expl\'icitamente en el
dise\~no de Spanner mientras que Kafka, empleado como backbone de mensajer\'ia as\'incrona, se
aproxima en esp\'iritu al extremo de disponibilidad que ilustra Dynamo.

% ============================================================
% ANALISIS DEL PFC
% ============================================================
\section{An\'alisis del modelo de consistencia del PFC}
\label{sec:analisis}

\subsection{Caracterizaci\'on del modelo implementado}
\label{subsec:caracterizacion}

El sistema del equipo ACC no implementa un \'unico modelo de consistencia uniforme, sino dos
regiones con garant\'ias distintas, separadas por la frontera as\'incrona de Kafka. La
Figura~\ref{fig:topologia} representa la topolog\'ia real desplegada.

\textbf{N\'ucleo transaccional (tickets, t\'ecnicos, autenticaci\'on).} El cl\'uster de CockroachDB
se despliega con tres nodos (\texttt{roach1}, \texttt{roach2}, \texttt{roach3}), cada uno anclado a
una \emph{locality} distinta (\texttt{zone=quevedo-centro}, \texttt{zone=quevedo-norte},
\texttt{zone=quevedo-sur}) mediante el flag \texttt{--locality} en el arranque de cada nodo. La
tabla principal \texttt{tickets} se particiona horizontalmente por lista (\texttt{PARTITION BY
LIST}) sobre la columna \texttt{zone}, con una partici\'on por zona operativa
(\texttt{tickets\_centro}, \texttt{tickets\_norte}, \texttt{tickets\_sur}) y una partici\'on por
defecto. Cada partici\'on se configura expl\'icitamente con \texttt{num\_replicas = 3} y una
restricci\'on de localidad preferente hacia su zona (\texttt{constraints =
\{"+zone=quevedo-centro": 1\}}, an\'aloga para las otras dos), de modo que el \emph{leaseholder}
de cada rango tiende a ubicarse en el nodo de su propia zona, pero las tres r\'eplicas Raft de cada
rango se mantienen sincronizadas entre los tres nodos del cl\'uster. Este es, por definici\'on,
un esquema de \textbf{replicaci\'on l\'ider--seguidor s\'incrona por rango}: cada escritura debe
alcanzar un qu\'orum Raft (dos de tres r\'eplicas) antes de confirmarse, y CockroachDB expone
\'unicamente aislamiento \texttt{SERIALIZABLE}, la garant\'ia m\'as fuerte del espectro descrito en
la Secci\'on~\ref{subsec:modelos} para operaciones transaccionales sobre el mismo rango. Tanto
\texttt{auth-service} como \texttt{ticket-service} se conectan a este cl\'uster con
\texttt{ddl-auto: none}, delegando el esquema -- incluida la partici\'on -- a los scripts de
\texttt{db-cluster/scripts}, y sin contrase\~na (modo \texttt{--insecure}, aceptable para el
entorno de desarrollo del PFC pero no para un despliegue productivo).

\textbf{Coreograf\'ia de clasificaci\'on (borde as\'incrono).} Cuando \texttt{ticket-service} crea
un ticket, publica un evento en el t\'opico \texttt{ticket.created} de Kafka. \texttt{ai-service}
(un servicio Python/FastAPI) consume ese t\'opico, ejecuta la clasificaci\'on y publica el
resultado en \texttt{ticket.classified}; \texttt{ticket-service} consume ese segundo t\'opico y
actualiza el ticket con la categor\'ia inferida. Este patr\'on -- una \emph{saga por
coreograf\'ia} sin coordinador central -- desacopla completamente ambos servicios: no existe
transacci\'on distribuida ni bloqueo entre la creaci\'on del ticket y su clasificaci\'on. La
consecuencia directa es que, entre el instante en que el ticket se crea (y ya es visible, sin
categor\'ia, para cualquier lector del n\'ucleo transaccional) y el instante en que
\texttt{ticket-service} aplica la actualizaci\'on proveniente de \texttt{ticket.classified}, existe
una ventana de \textbf{inconsistencia eventual} sobre el campo \texttt{category}: dos lectores
concurrentes del mismo ticket pueden observar temporalmente valores distintos de ese campo seg\'un
si su lectura ocurre antes o despu\'es de que la coreograf\'ia complete su ciclo. Kafka, en la
configuraci\'on de un solo bróker usada en el PFC, no aporta por s\'i mismo garant\'ias de
consistencia fuerte entre productor y consumidores; su contrato es de entrega ordenada por
partici\'on con procesamiento asincr\'onico, que es precisamente la propiedad que hace posible
desacoplar la disponibilidad del flujo principal de tickets respecto de la disponibilidad del
servicio de IA.

\begin{figure}[ht]
\centering
\resizebox{\linewidth}{!}{%
\begin{tikzpicture}[
  node distance=9mm and 14mm,
  box/.style={draw, rounded corners, align=center, minimum height=8mm, minimum width=24mm, font=\small},
  db/.style={draw, circle, align=center, minimum size=13mm, font=\scriptsize},
  strong/.style={-{Latex[length=2mm]}, thick},
  async/.style={-{Latex[length=2mm]}, thick, dashed},
  lbl/.style={font=\scriptsize, midway, fill=white, inner sep=1pt}
]
  % Capa de cliente y autenticacion
  \node[box] (front) {Frontend\\(cliente web)};
  \node[box, right=of front] (auth) {auth-service\\(JWT, :8001)};
  \node[box, right=of auth] (tix) {ticket-service\\(:8002)};

  % Cluster CockroachDB particionado por zona (region de consistencia fuerte)
  \node[db, below=16mm of front, xshift=6mm] (r1) {roach1\\quevedo\\centro};
  \node[db, right=of r1] (r2) {roach2\\quevedo\\norte};
  \node[db, right=of r2] (r3) {roach3\\quevedo\\sur};

  \node[draw, dashed, fit=(r1)(r2)(r3), inner sep=6mm, label={[font=\scriptsize]below:CockroachDB -- Raft, SERIALIZABLE, 3 replicas/rango}] (cluster) {};

  % Servicio de IA y Kafka (region de consistencia eventual)
  \node[box, right=of tix, yshift=-16mm] (ai) {ai-service\\(FastAPI, :8004)};
  \node[db, below=8mm of ai, minimum size=11mm] (mongo) {MongoDB};

  % Conexiones sincronas / fuertes
  \draw[strong] (front) -- node[lbl]{login} (auth);
  \draw[strong] (front) -- (tix);
  \draw[strong] (auth) -- node[lbl]{valida JWT} (tix);
  \draw[strong] (tix) -- node[lbl]{lecturas/escrituras SQL (qu\'orum Raft)} (cluster);

  % Conexiones asincronas / eventuales (Kafka)
  \draw[async] (tix.south) .. controls +(0,-10mm) and +(-10mm,0) .. node[lbl]{ticket.created} (ai.west);
  \draw[async] (ai.north) .. controls +(10mm,0) and +(0,-10mm) .. node[lbl]{ticket.classified} (tix.south east);
  \draw[strong] (ai) -- (mongo);

  \node[font=\scriptsize, align=left, below=2mm of ai, yshift=-16mm, xshift=6mm] (leyenda)
    {\textbf{Leyenda:}\\[1pt]
     \tikz\draw[strong,-{Latex[length=2mm]}] (0,0) -- (7mm,0); ~ s\'incrono / consistencia fuerte\\[1pt]
     \tikz\draw[async,-{Latex[length=2mm]}] (0,0) -- (7mm,0); ~ as\'incrono (Kafka) / consistencia eventual};
\end{tikzpicture}%
}
\caption{Topolog\'ia de replicaci\'on y consistencia del PFC del equipo ACC. El n\'ucleo
transaccional (izquierda, gris punteado) usa replicaci\'on s\'incrona por consenso Raft con
aislamiento serializable; la coreograf\'ia con \texttt{ai-service} (derecha) atraviesa la frontera
as\'incrona de Kafka y opera bajo consistencia eventual.}
\label{fig:topologia}
\end{figure}

\subsection{Justificaci\'on t\'ecnica ligada al dominio}
\label{subsec:justificacion}

La coexistencia de dos regiones de consistencia distintas no es una inconsistencia de dise\~no,
sino una lectura impl\'icita -- y correcta -- de PACELC aplicada a dos subproblemas del mismo
dominio con requisitos opuestos.

Para el n\'ucleo transaccional, el dominio impone invariantes que toleran mal la
consistencia relajada. La asignaci\'on de un t\'ecnico a un ticket, el c\'omputo del cumplimiento
de un SLA (\texttt{sla\_deadline}, \texttt{sla\_breached}) y el estado del ticket
(\texttt{status}) son datos cr\'iticos de negocio: una lectura obsoleta podr\'ia llevar a asignar
dos t\'ecnicos distintos al mismo ticket, o a subestimar un incumplimiento de SLA con
consecuencias contractuales frente al cliente. Bajo PACELC, el equipo ACC eligi\'o
\textbf{PC/EC}: tanto ante una partici\'on (P) como en operaci\'on normal (E), el sistema prioriza
consistencia (C) sobre disponibilidad/latencia, aceptando el costo de esperar el qu\'orum Raft en
cada escritura. Esta elecci\'on es coherente con el hecho de que CockroachDB solo ofrece
aislamiento serializable: el equipo no dise\~n\'o una relajaci\'on deliberada de consistencia para
el n\'ucleo, sino que adopt\'o la garant\'ia m\'as fuerte disponible en la plataforma elegida.

Para la coreograf\'ia de clasificaci\'on, en cambio, el dominio es tolerante a \emph{staleness}: la
categor\'ia autom\'atica de un ticket es un dato de enriquecimiento, \'util para priorizar la cola
de trabajo de los t\'ecnicos, pero su ausencia o retraso moment\'aneo no compromete la
correcci\'on del flujo de soporte -- el ticket sigue siendo v\'alido, visible y asignable sin
categor\'ia. Aqu\'i el equipo eligi\'o impl\'icitamente \textbf{PA/EL} para ese subproblema:
prioriza disponibilidad y baja latencia del flujo principal (\texttt{ticket-service} nunca espera
a \texttt{ai-service} para responder al cliente) sobre la consistencia inmediata del campo
\texttt{category}, exactamente el mismo compromiso que justifica el uso de consistencia eventual
en Dynamo~\cite{decandia2007}. La partici\'on geogr\'afica por zona (\texttt{PARTITION BY LIST} +
restricciones de localidad) responde adem\'as a un requisito de negocio real y verificable en el
c\'odigo: la mayor\'ia de las consultas y actualizaciones de un ticket las realiza el t\'ecnico de
su propia zona, por lo que anclar el \emph{leaseholder} de cada partici\'on al nodo local reduce la
latencia de las operaciones zonales m\'as frecuentes sin renunciar a las tres r\'eplicas
s\'incronas que garantizan tolerancia a la ca\'ida de cualquier nodo individual.

% ============================================================
% EVALUACION DE ALTERNATIVAS
% ============================================================
\section{Evaluaci\'on de alternativas}
\label{sec:alternativas}

Para contrastar el modelo actual se consideran dos alternativas de consistencia/replicaci\'on
razonables para el mismo dominio de gesti\'on de tickets.

\textbf{Alternativa 1 -- R\'eplica primario--secundario as\'incrona (estilo MongoDB replica
set).} Consistir\'ia en reemplazar el cl\'uster particionado de CockroachDB por una \'unica base de
datos primaria que replica de forma as\'incrona hacia una o m\'as secundarias, permitiendo lecturas
desde secundarias para descargar al primario. Es, de hecho, el mismo patr\'on que el propio PFC ya
usa -- pero para \texttt{ai-service}, no para el n\'ucleo -- con la instancia \texttt{mongo} del
\emph{stack} de mensajer\'ia, lo que la vuelve una alternativa concretamente evaluable en este
contexto.

\textbf{Alternativa 2 -- Almacenamiento sin l\'ider por qu\'orum (estilo Dynamo/Cassandra).}
Consistir\'ia en distribuir la tabla \texttt{tickets} sobre un anillo de nodos sin l\'ider fijo,
con factor de r\'eplicaci\'on $N=3$ y par\'ametros de qu\'orum de lectura/escritura ajustables
($R+W>N$ para consistencia fuerte ajustada, o $R+W\le N$ para favorecer disponibilidad y
latencia), resolviendo conflictos de escrituras concurrentes mediante vectores de versi\'on o
\emph{last-write-wins}.

La Tabla~\ref{tab:comparativa} contrasta el modelo actual frente a ambas alternativas seg\'un cinco
criterios t\'ecnicos.

\begin{table}[ht]
\centering
\small
\caption{Comparaci\'on del modelo actual del PFC (ACC) frente a dos alternativas de
consistencia/replicaci\'on.}
\label{tab:comparativa}
\begin{tabular}{@{}p{2.6cm}p{3.7cm}p{3.7cm}p{3.7cm}@{}}
\toprule
\textbf{Criterio} & \textbf{Modelo actual}\newline(CockroachDB Raft + Kafka) & \textbf{Alt. 1}\newline(primario--secundario as\'incrono) & \textbf{Alt. 2}\newline(sin l\'ider por qu\'orum) \\
\midrule
Consistencia &
Serializable (fuerte) en el n\'ucleo; eventual solo en el campo de clasificaci\'on &
Fuerte en el primario; eventual en lecturas desde secundarias &
Ajustable v\'ia $R,W$; eventual por defecto, fuerte solo con $R+W>N$ \\
\addlinespace
Disponibilidad ante ca\'ida de nodo &
Alta: tolera la ca\'ida de 1 de 3 nodos sin p\'erdida de qu\'orum Ra. &
Baja para escritura si cae el primario (requiere failover); alta para lectura &
Muy alta: cualquier nodo disponible puede servir la operaci\'on \\
\addlinespace
Latencia de escritura &
Media: espera qu\'orum Raft (2 de 3), atenuada por afinidad de zona &
Baja (el primario confirma sin esperar secundarias) &
Baja a media, seg\'un $W$ configurado \\
\addlinespace
Tolerancia a particiones &
Alta dentro del cl\'uster; el n\'ucleo prioriza C sobre A ante partici\'on &
El primario aislado sigue escribiendo; secundarias divergen &
Alta: dise\~nado expl\'icitamente para operar particionado (P) \\
\addlinespace
Complejidad operativa &
Media: partici\'on y localidades declarativas, sin resoluci\'on manual de conflictos &
Baja: patr\'on est\'andar y bien soportado &
Alta: exige gesti\'on expl\'icita de conflictos y de qu\'orum \\
\bottomrule
\end{tabular}
\end{table}

Frente a la Alternativa 1, el modelo actual ofrece una garant\'ia estrictamente m\'as fuerte para
el n\'ucleo (serializable multi-nodo frente a fuerte-solo-en-el-primario), al costo de una latencia
de escritura algo mayor; dado que el dominio penaliza m\'as una asignaci\'on incorrecta de t\'ecnico
que unos pocos milisegundos adicionales por transacci\'on, el modelo actual domina para el
n\'ucleo. Frente a la Alternativa 2, el modelo actual sacrifica algo de disponibilidad y de
tolerancia a particiones extremas a cambio de eliminar por completo la complejidad de resoluci\'on
de conflictos concurrentes, complejidad que un equipo de PFC de alcance acad\'emico dif\'icilmente
podr\'ia operar de forma confiable dentro del tiempo del curso.

% ============================================================
% DISCUSION CRITICA
% ============================================================
\section{Discusi\'on cr\'itica}
\label{sec:discusion}

¿Es el modelo de consistencia elegido el \'optimo para el dominio del sistema? La respuesta
razonada es: \textbf{s\'i, para el n\'ucleo transaccional, con reservas operativas}, y
\textbf{s\'i, de forma casi \'optima, para la coreograf\'ia de clasificaci\'on}.

Para el n\'ucleo, la elecci\'on de consistencia serializable v\'ia Raft es defendible bajo PACELC
porque el dominio -- asignaci\'on de t\'ecnicos y c\'omputo de SLA -- es exactamente el tipo de
carga para la que CAP/PACELC recomienda priorizar C sobre A/L: los costos de una violaci\'on de
consistencia (doble asignaci\'on, SLA mal calculado) son altos y dif\'icilmente reversibles,
mientras que el costo de la latencia adicional del qu\'orum Raft es acotado y predecible dentro de
un cl\'uster de tres nodos en la misma regi\'on. Sin embargo, subsisten dos limitaciones concretas.
Primera, el cl\'uster corre en modo \texttt{--insecure} y sin r\'eplicas fuera de la m\'aquina de
desarrollo: la tolerancia a fallos que ofrece Raft es real frente a la ca\'ida de un proceso, pero
no frente a un fallo del host completo, lo que en producci\'on exigir\'ia nodos en m\'aquinas o
zonas de disponibilidad f\'isicamente independientes. Segunda, anclar el \emph{leaseholder} de
cada partici\'on a su zona reduce la latencia de las operaciones locales, pero cualquier consulta
que cruce zonas (por ejemplo, un reporte agregado de SLA a nivel de ciudad) sigue pagando el costo
de un qu\'orum Raft con al menos una r\'eplica remota, un compromiso que el equipo no documenta
expl\'icitamente y que conviene medir.

Para la coreograf\'ia de clasificaci\'on, la consistencia eventual v\'ia Kafka es la elecci\'on
correcta en esp\'iritu -- el enriquecimiento no debe bloquear el flujo cr\'itico -- pero el PFC no
implementa ninguna de las mitigaciones habituales para acotar el impacto de esa eventualidad: no
hay un l\'imite expl\'icito de tiempo m\'aximo de convergencia (\emph{staleness bound}), ni una
se\~nal visible en la interfaz de que la categor\'ia de un ticket est\'a "pendiente de
clasificar", ni una pol\'itica de reintentos o \emph{dead-letter queue} documentada si
\texttt{ai-service} cae mientras hay mensajes en \texttt{ticket.created}. Estas mitigaciones no
cambiar\'ian el modelo de consistencia elegido, pero s\'i reducir\'ian el riesgo de que la
consistencia eventual se perciba, err\'oneamente, como un dato perdido en vez de un dato en
tr\'ansito. En conjunto, el modelo h\'ibrido implementado es una lectura correcta de CAP/PACELC
aplicada de forma diferenciada a dos subproblemas con requisitos opuestos dentro del mismo
dominio, y su principal margen de mejora no est\'a en cambiar el modelo de consistencia, sino en
hacer expl\'icitas y observables las garant\'ias -- y las ausencias de garant\'ia -- que ya
implementa.

% ============================================================
% CONCLUSIONES
% ============================================================
\section{Conclusiones}

El PFC del equipo ACC implementa, de facto, un modelo de consistencia h\'ibrido: consistencia
fuerte (serializable, v\'ia consenso Raft con replicaci\'on s\'incrona de tres r\'eplicas por
partici\'on geogr\'afica) para el n\'ucleo transaccional de tickets, y consistencia eventual (v\'ia
una coreograf\'ia as\'incrona sobre Kafka) para el enriquecimiento de clasificaci\'on autom\'atica.
Esta combinaci\'on no fue necesariamente el resultado de un an\'alisis expl\'icito de CAP/PACELC
por parte del equipo, pero es, retrospectivamente, coherente con lo que ambos marcos recomiendan
para un dominio donde coexisten invariantes cr\'iticos de negocio (asignaci\'on, SLA) con tareas de
enriquecimiento tolerantes a retraso (clasificaci\'on). La comparaci\'on con dos alternativas de
dise\~no -- r\'eplica primario--secundario as\'incrona y almacenamiento sin l\'ider por qu\'orum --
confirma que ninguna domina estrictamente al modelo actual para el n\'ucleo transaccional, y que la
complejidad adicional de un esquema sin l\'ider no se justifica en el alcance de un PFC
acad\'emico. El aprendizaje principal para el equipo es que la optimalidad del modelo actual es
condicional: se sostiene mientras el dominio siga privilegiando correcci\'on sobre latencia en el
n\'ucleo y tolere retraso en el enriquecimiento, y se debilita si el sistema creciera hacia
m\'ultiples regiones geogr\'aficas reales o si el negocio empezara a exigir garant\'ias de tiempo
de convergencia sobre la clasificaci\'on autom\'atica.

% ============================================================
% DECLARACION DE USO DE IA
% ============================================================
\section{Declaraci\'on de uso de IA generativa}

Para la elaboraci\'on de este informe se utiliz\'o Claude (Anthropic), un asistente de IA
generativa, como herramienta de apoyo en las siguientes tareas: (i) inspecci\'on del c\'odigo
fuente y de los archivos de configuraci\'on del PFC del equipo ACC (\texttt{docker-compose},
\texttt{application.yml}, scripts SQL de partici\'on y zonas) para caracterizar con precisi\'on el
modelo de consistencia y la estrategia de replicaci\'on realmente implementados; (ii) organizaci\'on
del documento conforme a la estructura obligatoria de la gu\'ia TA-IND-02; y (iii) redacci\'on
inicial del marco te\'orico, del an\'alisis del modelo del PFC, de la tabla comparativa y de la
discusi\'on cr\'itica. El estudiante revis\'o el contenido generado, verific\'o su correspondencia
con el c\'odigo real del PFC y asumi\'o la postura cr\'itica final presentada en la
Secci\'on~\ref{sec:discusion}.

% ============================================================
% REFERENCIAS (IEEE, manuales, sin bibtex)
% ============================================================
\begin{thebibliography}{6}

\bibitem{abadi2012}
D.~J. Abadi, ``Consistency Tradeoffs in Modern Distributed Database System Design: CAP is Only
Part of the Story,'' \emph{Computer}, vol.~45, no.~2, pp.~37--42, 2012. doi:
10.1109/MC.2012.33

\bibitem{corbett2013}
J.~C. Corbett \emph{et al.}, ``Spanner: Google's Globally Distributed Database,'' \emph{ACM
Transactions on Computer Systems}, vol.~31, no.~3, pp.~1--22, 2013, art.\ no.~8. doi:
10.1145/2491245

\bibitem{decandia2007}
G.~DeCandia \emph{et al.}, ``Dynamo: Amazon's Highly Available Key-value Store,'' in
\emph{Proceedings of the 21st ACM SIGOPS Symposium on Operating Systems Principles (SOSP
'07)}, 2007, pp.~205--220. doi: 10.1145/1294261.1294281

\bibitem{gilbert2002}
S.~Gilbert and N.~Lynch, ``Brewer's Conjecture and the Feasibility of Consistent, Available,
Partition-Tolerant Web Services,'' \emph{ACM SIGACT News}, vol.~33, no.~2, pp.~51--59, 2002.
doi: 10.1145/564585.564601

\bibitem{vansteen2017}
M.~van Steen and A.~S. Tanenbaum, \emph{Distributed Systems}, 3rd ed. distributed-systems.net,
2017. isbn: 978-1543057386.

\bibitem{viotti2016}
P.~Viotti and M.~Vukoli\'c, ``Consistency in Non-Transactional Distributed Storage Systems,''
\emph{ACM Computing Surveys}, vol.~49, no.~1, pp.~1--34, 2016, art.\ no.~19. doi:
10.1145/2926965

\end{thebibliography}

\end{document}
